\section{Conclusion: geometry as source code (under strict layering)}
\label{sec:conclusion}

Under the strict layering discipline of Section~\ref{sec:layers_axioms}, the \emph{Ramanujan holographic scanning principle} may be summarized as three constitution clauses:
\begin{enumerate}
  \item \textbf{Time} is the iteration count of an intrinsic unitary scan (O1, O3, O6).
  \item \textbf{Probability} is an induced measure produced by finite-resolution projection readout (O5), not an external sampling postulate.
  \item \textbf{Discreteness} is a cusp interface: $q$-expansions produce arithmetic coefficient data, constrained by Hecke/prime structure and encoded canonically via Ostrowski/Zeckendorf (Sections~\ref{sec:readout_as_q_expansion}--\ref{sec:canonical_coding}).
\end{enumerate}

In this sense, HPA is not an auxiliary coding trick but the minimal arithmetic implementation of modular geometry under finite-resolution readout, and the $\Omega$ framework supplies an ontological container compatible with finite information and holographic encoding. The accompanying toy experiments provide reproducible evidence that the proposed chain is mathematically coherent and computationally checkable.


